% Book 1 appendix: typing drills. Owned by g3/01 (typing-with-ten-fingers), which teaches the technique.
% Provenance (never printed): see sources/book1/BRIEFS.md, section 8.
% Every \omdrill line carries a source comment "% drill: <layout> keys=<keys> [shift]": the xkb keys
% (names or ranges, or "all") whose level-1 symbols, and level-2 symbols with "shift", the line may use;
% the space bar is always allowed. Every \ompassage count is the length of its text. Both are checked by
% code/typing/01-typing-drills/tests/test_drills.py against the system layout files (us, fr, de).
% The words are English in every layout; a language edition may use its own words.
% ledger: typ:layout-us, typ:layout-fr, typ:layout-de
\chapter{Typing Drills}\label{ch:typ:typing-drills}

This appendix holds practice lines for \omterm{def:g3:typing-with-ten-fingers:touch-typing}{touch typing}, for three \omterm{def:g3:typing-with-ten-fingers:layout}{keyboard
layouts}: QWERTY, AZERTY and QWERTZ. The technique (the \omterm{def:g3:typing-with-ten-fingers:home-row}{home row}, which finger
strikes which key, \key{Shift} with the other hand) is taught in
\cref{ch:g3:typing-with-ten-fingers}. Find the layout of your keyboard by
reading the first six letters of its top row of letters, and use only its
section.

\section{How to use the drills}

\begin{method}[A practice session]\label{met:typ:typing-drills:session}
Practise ten minutes a day rather than one long hour a week.
\begin{enumerate}
\item Sit well, put your index fingers on the bumps of F and J, and let the
  other fingers fall onto the \omterm{def:g3:typing-with-ten-fingers:home-row}{home row}.
\item Warm up with the first lines of the \omterm{def:g3:typing-with-ten-fingers:home-row}{home row}, slowly.
\item Take the next lines in order. Type each line three times without looking
  at the keys. If you make more than two mistakes in a line, type it again,
  more slowly.
\item Move on to the next group of lines only when the last one feels easy.
\item Finish with one timed passage.
\end{enumerate}
\end{method}

To time a passage, ask somebody to watch a clock, or use a timer. Type the
passage once, then divide its number of \omterm{def:g7:strings:string}{characters} by five to get words, and
divide by the number of minutes: this gives your speed in \omterm{def:g3:typing-with-ten-fingers:wpm}{words per minute}
(\cref{def:g3:typing-with-ten-fingers:wpm}). Write down the speed and the
number of mistakes, with the date, and watch both improve over the weeks.
Accuracy comes first: a fast line full of mistakes is not progress.

% ---------------------------------------------------------------------------
\section{QWERTY}

\begin{omfigure}
\omkeyboard[fingers, home]{us}

{\small The QWERTY keyboard, with the zone of each finger: little fingers
mint, ring fingers yellow, middle fingers red, index fingers blue, thumbs
grey.}
\label{fig:typ:typing-drills:us}
\end{omfigure}

\subsection*{Home row}

\omdrill{fff jjj fff jjj fjf jfj fjf jfj}% drill: us keys=AC04,AC07
\omdrill{ddd kkk ddd kkk dkd kdk dkd kdk}% drill: us keys=AC03,AC08
\omdrill{sss lll sss lll sls lsl sls lsl}% drill: us keys=AC02,AC09
\omdrill{aaa ;;; aaa ;;; a;a ;a; a;a ;a;}% drill: us keys=AC01,AC10
\omdrill{asdf jkl; asdf jkl; fdsa ;lkj fdsa ;lkj}% drill: us keys=AC01-AC10
\omdrill{fgf jhj fgf jhj gf hj gf hj fg jh}% drill: us keys=AC04-AC07
\omdrill{a lad; a lass; a flask; a hall; a flag;}% drill: us keys=AC01-AC10
\omdrill{ask dad; add salad; all fall; shall;}% drill: us keys=AC01-AC10
\omdrill{all lads had a glass; dad has a flask;}% drill: us keys=AC01-AC10

\subsection*{Top row}

\omdrill{frf juj frf juj ftf jyj ftf jyj}% drill: us keys=AC04,AC07,AD04-AD07
\omdrill{ded kik ded kik dek kid dek kid}% drill: us keys=AC03,AC08,AD03,AD08
\omdrill{sws lol sws lol swl los swl los}% drill: us keys=AC02,AC09,AD02,AD09
\omdrill{aqa ;p; aqa ;p; qa; pa; qa; pa;}% drill: us keys=AC01,AC10,AD01,AD10
\omdrill{the red fish; a quiet path; it is our street;}% drill: us keys=AC01-AC10,AD01-AD10
\omdrill{where were you; type faster; we said it;}% drill: us keys=AC01-AC10,AD01-AD10

\subsection*{Bottom row}

\omdrill{fvf jmj fvf jmj fbf jnj fbf jnj}% drill: us keys=AC04,AC07,AB04-AB07
\omdrill{dcd k,k dcd k,k dc, kc, dc, kc,}% drill: us keys=AC03,AC08,AB03,AB08
\omdrill{sxs l.l sxs l.l sx. lx. sx. lx.}% drill: us keys=AC02,AC09,AB02,AB09
\omdrill{aza ;/; aza ;/; za/ az; za/ az;}% drill: us keys=AC01,AC10,AB01,AB10
\omdrill{my box, a calm zebra, five big jams.}% drill: us keys=AC01-AC10,AD01-AD10,AB01-AB10
\omdrill{the quick brown fox jumps over the lazy dog.}% drill: us keys=AC01-AC10,AD01-AD10,AB01-AB10

\subsection*{Capitals}

\omdrill{Ada Fred Kim Lola Jack Sam Hugo Rita}% drill: us keys=AC01-AC10,AD01-AD10,AB01-AB10 shift
\omdrill{Paris London Oslo Tokyo Cairo Lima}% drill: us keys=AC01-AC10,AD01-AD10,AB01-AB10 shift

\subsection*{Digits and signs}

\omdrill{1 2 3 4 5 6 7 8 9 0 10 20 30 40 50}% drill: us keys=AE01-AE10
\omdrill{12 34 56 78 90 2026 1874 365 24 60}% drill: us keys=AE01-AE10
\omdrill{Yes. No? Wait! one, two; three: four - five (six)}% drill: us keys=all shift

\subsection*{Sentences}

\omdrill{Kim has 3 cats and 2 dogs. Are they all asleep?}% drill: us keys=all shift
\omdrill{We meet at 9, then walk 4 km to the lake.}% drill: us keys=all shift

\subsection*{Timed passages}

\ompassage{83}{The cat sat on the warm mat. A small dog ran past, and the cat did not move at all.}% drill: us keys=all shift
\ompassage{207}{Every morning, Lina walks to school with her brother. They cross the park, wave to the baker and count the birds on the old wall. Yesterday they saw 14 sparrows and 2 crows. Today it is raining, so they run.}% drill: us keys=all shift
\ompassage{375}{A computer follows instructions, one after another, very fast. It keeps the work in progress in its memory and the saved work in its storage. When you type, each key sends a small message to the processor, which puts a letter on the screen. A good typist keeps the eyes on the text and the fingers on the home row, and types slowly at first. Speed comes later, with practice.}% drill: us keys=all shift

% ---------------------------------------------------------------------------
\section{AZERTY}

\begin{omfigure}
\omkeyboard[fingers, home]{fr}

{\small The AZERTY keyboard, with the same finger zones. A and Q, Z and W are
swapped compared with QWERTY, M sits on the \omterm{def:g3:typing-with-ten-fingers:home-row}{home row}, and the digits of the
top row are typed with \key{Shift}.}
\label{fig:typ:typing-drills:fr}
\end{omfigure}

\subsection*{Home row}

The \omterm{def:g3:typing-with-ten-fingers:home-row}{home row} of AZERTY has no vowel, so its first lines are letters only;
words come with the top row.

\omdrill{fff jjj fff jjj fjf jfj fjf jfj}% drill: fr keys=AC04,AC07
\omdrill{ddd kkk ddd kkk dkd kdk dkd kdk}% drill: fr keys=AC03,AC08
\omdrill{sss lll sss lll sls lsl sls lsl}% drill: fr keys=AC02,AC09
\omdrill{qqq mmm qqq mmm qmq mqm qmq mqm}% drill: fr keys=AC01,AC10
\omdrill{qsdf jklm qsdf jklm fdsq mlkj fdsq mlkj}% drill: fr keys=AC01-AC10
\omdrill{fgf jhj fgf jhj gf hj gf hj fg jh}% drill: fr keys=AC04-AC07

\subsection*{Top row}

\omdrill{frf juj frf juj ftf jyj ftf jyj}% drill: fr keys=AC04,AC07,AD04-AD07
\omdrill{ded kik ded kik dek kid dek kid}% drill: fr keys=AC03,AC08,AD03,AD08
\omdrill{szs lol szs lol szl los szl los}% drill: fr keys=AC02,AC09,AD02,AD09
\omdrill{qaq mpm qaq mpm aqm pmq aqm pmq}% drill: fr keys=AC01,AC10,AD01,AD10
\omdrill{the red fish a quiet path my mother is here}% drill: fr keys=AC01-AC10,AD01-AD10
\omdrill{make a list ask her type faster it is our street}% drill: fr keys=AC01-AC10,AD01-AD10

\subsection*{Bottom row}

\omdrill{fvf jnj fvf jnj fbf j,j fbf j,j}% drill: fr keys=AC04,AC07,AB04-AB07
\omdrill{dcd k;k dcd k;k dc; kc; dc; kc;}% drill: fr keys=AC03,AC08,AB03,AB08
\omdrill{sxs l:l sxs l:l sx: lx: sx: lx:}% drill: fr keys=AC02,AC09,AB02,AB09
\omdrill{qwq m!m qwq m!m wq! mw! wq! mw!}% drill: fr keys=AC01,AC10,AB01,AB10
\omdrill{my box, a calm zebra, five big jams}% drill: fr keys=AC01-AC10,AD01-AD10,AB01-AB10
\omdrill{the quick brown fox jumps over the lazy dog}% drill: fr keys=AC01-AC10,AD01-AD10,AB01-AB10

\subsection*{Capitals}

\omdrill{Ada Fred Kim Lola Jack Sam Hugo Rita}% drill: fr keys=AC01-AC10,AD01-AD10,AB01-AB10 shift
\omdrill{Paris London Oslo Tokyo Cairo Lima}% drill: fr keys=AC01-AC10,AD01-AD10,AB01-AB10 shift

\subsection*{Digits, accented letters and signs}

On AZERTY, the digits need \key{Shift}; without it, the same keys type
accented letters and signs.

\omdrill{1 2 3 4 5 6 7 8 9 0 10 20 30 40 50}% drill: fr keys=AE01-AE10 shift
\omdrill{12 34 56 78 90 2026 1874 365 24 60}% drill: fr keys=AE01-AE10 shift
\omdrill{é è à ç ù é è à ç ù éà èç ùé}% drill: fr keys=AE02,AE07,AE09,AE10,AC11
\omdrill{café déjà vu à la carte façade}% drill: fr keys=all
\omdrill{Yes. No? Wait! one, two; three: four - five (six)}% drill: fr keys=all shift

\subsection*{Sentences}

\omdrill{Kim has 3 cats and 2 dogs. Are they all asleep?}% drill: fr keys=all shift
\omdrill{We meet at 9, then walk 4 km to the lake.}% drill: fr keys=all shift

\subsection*{Timed passages}

\ompassage{83}{The cat sat on the warm mat. A small dog ran past, and the cat did not move at all.}% drill: fr keys=all shift
\ompassage{207}{Every morning, Lina walks to school with her brother. They cross the park, wave to the baker and count the birds on the old wall. Yesterday they saw 14 sparrows and 2 crows. Today it is raining, so they run.}% drill: fr keys=all shift
\ompassage{375}{A computer follows instructions, one after another, very fast. It keeps the work in progress in its memory and the saved work in its storage. When you type, each key sends a small message to the processor, which puts a letter on the screen. A good typist keeps the eyes on the text and the fingers on the home row, and types slowly at first. Speed comes later, with practice.}% drill: fr keys=all shift

% ---------------------------------------------------------------------------
\section{QWERTZ}

\begin{omfigure}
\omkeyboard[fingers, home]{de}

{\small The QWERTZ keyboard, with the same finger zones. Z and Y are swapped
compared with QWERTY, and the right little finger reaches Ö, Ä and Ü.}
\label{fig:typ:typing-drills:de}
\end{omfigure}

\subsection*{Home row}

\omdrill{fff jjj fff jjj fjf jfj fjf jfj}% drill: de keys=AC04,AC07
\omdrill{ddd kkk ddd kkk dkd kdk dkd kdk}% drill: de keys=AC03,AC08
\omdrill{sss lll sss lll sls lsl sls lsl}% drill: de keys=AC02,AC09
\omdrill{aaa ööö aaa ööö aöa öaö aöa öaö}% drill: de keys=AC01,AC10
\omdrill{asdf jklö asdf jklö fdsa ölkj fdsa ölkj}% drill: de keys=AC01-AC10
\omdrill{fgf jhj fgf jhj gf hj gf hj fg jh}% drill: de keys=AC04-AC07
\omdrill{a lad a lass a flask a hall a flag}% drill: de keys=AC01-AC10
\omdrill{ask dad add salad all fall shall}% drill: de keys=AC01-AC10
\omdrill{all lads had a glass dad has a flask}% drill: de keys=AC01-AC10

\subsection*{Top row}

\omdrill{frf juj frf juj ftf jzj ftf jzj}% drill: de keys=AC04,AC07,AD04-AD07
\omdrill{ded kik ded kik dek kid dek kid}% drill: de keys=AC03,AC08,AD03,AD08
\omdrill{sws lol sws lol swl los swl los}% drill: de keys=AC02,AC09,AD02,AD09
\omdrill{aqa öpö aqa öpö qaö pöa qaö pöa}% drill: de keys=AC01,AC10,AD01,AD10
\omdrill{the red fish a quiet path it is our street}% drill: de keys=AC01-AC10,AD01-AD10
\omdrill{zero to three we said it a fast tiger}% drill: de keys=AC01-AC10,AD01-AD10

\subsection*{Bottom row}

\omdrill{fvf jmj fvf jmj fbf jnj fbf jnj}% drill: de keys=AC04,AC07,AB04-AB07
\omdrill{dcd k,k dcd k,k dc, kc, dc, kc,}% drill: de keys=AC03,AC08,AB03,AB08
\omdrill{sxs l.l sxs l.l sx. lx. sx. lx.}% drill: de keys=AC02,AC09,AB02,AB09
\omdrill{aya ö-ö aya ö-ö ya- ay- ya- ay-}% drill: de keys=AC01,AC10,AB01,AB10
\omdrill{my box, a calm zebra, five big jams.}% drill: de keys=AC01-AC10,AD01-AD10,AB01-AB10
\omdrill{the quick brown fox jumps over the lazy dog.}% drill: de keys=AC01-AC10,AD01-AD10,AB01-AB10

\subsection*{Capitals}

\omdrill{Ada Fred Kim Lola Jack Sam Hugo Rita}% drill: de keys=AC01-AC10,AD01-AD10,AB01-AB10 shift
\omdrill{Paris London Oslo Tokyo Cairo Lima}% drill: de keys=AC01-AC10,AD01-AD10,AB01-AB10 shift

\subsection*{Digits, letters with dots and signs}

\omdrill{1 2 3 4 5 6 7 8 9 0 10 20 30 40 50}% drill: de keys=AE01-AE10
\omdrill{12 34 56 78 90 2026 1874 365 24 60}% drill: de keys=AE01-AE10
\omdrill{ä ö ü ß äöü ßüöä üß äö}% drill: de keys=AC10,AC11,AD11,AE11
\omdrill{Yes. No? Wait! one, two; three: four - five (six)}% drill: de keys=all shift

\subsection*{Sentences}

\omdrill{Kim has 3 cats and 2 dogs. Are they all asleep?}% drill: de keys=all shift
\omdrill{We meet at 9, then walk 4 km to the lake.}% drill: de keys=all shift

\subsection*{Timed passages}

\ompassage{83}{The cat sat on the warm mat. A small dog ran past, and the cat did not move at all.}% drill: de keys=all shift
\ompassage{207}{Every morning, Lina walks to school with her brother. They cross the park, wave to the baker and count the birds on the old wall. Yesterday they saw 14 sparrows and 2 crows. Today it is raining, so they run.}% drill: de keys=all shift
\ompassage{375}{A computer follows instructions, one after another, very fast. It keeps the work in progress in its memory and the saved work in its storage. When you type, each key sends a small message to the processor, which puts a letter on the screen. A good typist keeps the eyes on the text and the fingers on the home row, and types slowly at first. Speed comes later, with practice.}% drill: de keys=all shift
